\apendice{Anexo de sostenibilización curricular}

\section{Introducción}

En este anexo se presenta una reflexión sobre los aspectos de sostenibilidad relacionados con el desarrollo del Trabajo de Fin de Grado. En particular, se analizan las implicaciones del proyecto desde el punto de vista tecnológico, educativo y social, en línea con las competencias de sostenibilidad promovidas por la CRUE.

El sistema desarrollado, basado en la recopilación y análisis de datos durante el uso de Blender, permite abordar la sostenibilidad desde una perspectiva digital, centrada en la eficiencia, la optimización de procesos y la mejora del aprendizaje.

\section{Sostenibilidad tecnológica}

Desde el punto de vista tecnológico, el proyecto promueve la eficiencia en el uso de recursos. El sistema ha sido diseñado para registrar únicamente cambios relevantes en la interacción del usuario, evitando la generación de datos redundantes. Esta decisión reduce el volumen de información almacenada y, por tanto, el consumo de recursos de almacenamiento y procesamiento.

Además, el uso de archivos CSV como formato de persistencia contribuye a la simplicidad y ligereza del sistema, evitando dependencias complejas o infraestructuras costosas. Este enfoque permite minimizar el impacto energético asociado al procesamiento de datos.

Asimismo, el desarrollo se basa en herramientas de software libre como Blender y Python, lo que favorece la sostenibilidad tecnológica al promover soluciones accesibles, reutilizables y mantenibles en el tiempo.

\section{Sostenibilidad educativa}

El proyecto tiene una clara aplicación en el ámbito educativo, donde permite analizar el proceso de aprendizaje de los usuarios, más allá del resultado final. Mediante la recopilación de datos de interacción, es posible estudiar cómo los estudiantes trabajan en entornos de modelado 3D, identificar patrones de comportamiento y detectar posibles dificultades.

Este enfoque contribuye a una educación más sostenible, ya que permite mejorar los procesos de enseñanza y aprendizaje, adaptándolos a las necesidades reales de los estudiantes. En lugar de centrarse únicamente en la evaluación del resultado final, el sistema proporciona herramientas para comprender el proceso, lo que favorece un aprendizaje más eficiente y personalizado.

\section{Sostenibilidad social}

Desde el punto de vista social, el uso de software libre garantiza la accesibilidad del sistema a un amplio número de usuarios, sin barreras económicas. Blender y Python son herramientas ampliamente disponibles, lo que permite que el sistema desarrollado pueda ser utilizado en diferentes contextos educativos y profesionales.

Además, el sistema incorpora mecanismos de consentimiento para la recopilación de datos, respetando la privacidad del usuario y cumpliendo con la normativa vigente. Este aspecto refuerza el compromiso ético del proyecto, garantizando un uso responsable de la información.

\section{Conclusión}

En conjunto, el proyecto contribuye a la sostenibilidad desde múltiples perspectivas. Por un lado, promueve la eficiencia tecnológica mediante la optimización del registro y procesamiento de datos. Por otro, aporta valor en el ámbito educativo al mejorar la comprensión del proceso de aprendizaje. Finalmente, garantiza la accesibilidad y el respeto a la privacidad, alineándose con principios de sostenibilidad social.

Estas características reflejan la integración de competencias de sostenibilidad en el desarrollo del Trabajo de Fin de Grado, demostrando una visión global del impacto de la tecnología en la sociedad.